% © 2026 Ing. David Jaroš — CC BY-NC-ND 4.0
%
% This work is licensed under a Creative Commons Attribution-NonCommercial-NoDerivatives
% 4.0 International License (CC BY-NC-ND 4.0).
%
% License History: Earlier drafts (up to v0.3) were released under CC BY 4.0.
% From v0.4 onward, all material is released under CC BY-NC-ND 4.0 to protect
% the integrity of the theoretical work during ongoing academic development.
%
% See LICENSE.md for full license text.
%
% File: canonical/alpha/veff_corrected.tex
% Purpose: First-principles derivation of V_eff(n) from the UBT action,
%          correct stationarity analysis, and classification of the B-coefficient.
%          Deliverable from the Alpha Critical Fix task (2026-05-02).
%          Supersedes erroneous formulas in appendix_alpha_geometry.tex (pre-2026-05-02).
% Related: canonical/alpha/veff_corrected_statement.tex  (formal theorem list)
%          reports/veff_extremum_audit.md                (derivative check table)
%          reports/alpha_consistency_verdict.md          (route verdict)

\ifdefined\INCLUDEMODE
\else
  \documentclass[12pt,a4paper]{article}
  \usepackage[a4paper, margin=2.5cm]{geometry}
  \usepackage{amsmath, amssymb, amsthm}
  \usepackage{mathtools}
  \usepackage{hyperref}
  \usepackage{xcolor}
  \usepackage{booktabs}
  \usepackage{longtable}
  \usepackage{mdframed}
  \usepackage{tcolorbox}
  \tcbuselibrary{skins,breakable}
  \usepackage{array}

  \newtheorem{theorem}{Theorem}[section]
  \newtheorem{lemma}[theorem]{Lemma}
  \newtheorem{proposition}[theorem]{Proposition}
  \newtheorem{corollary}[theorem]{Corollary}
  \theoremstyle{definition}
  \newtheorem{definition}[theorem]{Definition}
  \theoremstyle{remark}
  \newtheorem{remark}[theorem]{Remark}

  \newmdenv[backgroundcolor=yellow!20, linecolor=orange, linewidth=1.5pt]{gapbox}
  \newmdenv[backgroundcolor=green!15, linecolor=green!60!black, linewidth=1.5pt]{resultbox}
  \newmdenv[backgroundcolor=red!8,   linecolor=red!60,         linewidth=1.5pt]{warnbox}
  \newmdenv[backgroundcolor=blue!6,  linecolor=blue!40,        linewidth=1.5pt]{infobox}

  \newcommand{\BB}{\mathbb{B}}
  \newcommand{\CC}{\mathbb{C}}
  \newcommand{\RR}{\mathbb{R}}
  \newcommand{\HH}{\mathbb{H}}
  \newcommand{\ZZ}{\mathbb{Z}}
  \newcommand{\NN}{\mathbb{N}}
  \newcommand{\Tr}{\mathrm{Tr}}

  \newcommand{\PROVED}{\textbf{[Proved]}}
  \newcommand{\OPEN}{\textbf{[Open]}}
  \newcommand{\COND}{\textbf{[COND]}}
  \newcommand{\MC}{\textbf{[MC]}}
  \newcommand{\Lone}{\textbf{[L1]}}
  \newcommand{\Lzero}{\textbf{[L0]}}
  \newcommand{\PHENOM}{\textbf{[PHENOM]}}

  \title{\textbf{First-Principles Derivation of $V_{\mathrm{eff}}(n)$\\
         and the $n^* = 137$ Stationarity Analysis}\\
         \large Alpha Critical Fix — Deliverable}
  \author{Ing.\ David Jaro\v{s}}
  \date{May 2026}

  \begin{document}
  \maketitle
  \tableofcontents
  \newpage
\fi

%──────────────────────────────────────────────────────────────────────────────
\section{Purpose}
\label{sec:vc2:purpose}
%──────────────────────────────────────────────────────────────────────────────

This document provides the \emph{first-principles derivation} of the effective
potential $V_{\mathrm{eff}}(n)$ from the UBT action $S[\Theta]$, derives its
stationarity condition without assumptions, and classifies all $B$-coefficient
values found in the repository.

\begin{warnbox}
\textbf{Hard rules (maintained throughout)}:
\begin{itemize}
  \item No functional form of $V_{\mathrm{eff}}$ is assumed — it is derived.
  \item No constant is chosen to reproduce $\alpha$ or $137$.
  \item Every equation is either derived or explicitly labelled with its
        proof status~(\PROVED, \COND, \MC, \PHENOM, \OPEN).
\end{itemize}
\end{warnbox}

%──────────────────────────────────────────────────────────────────────────────
\section{Ingredients: KK Spectrum and Entropy Term}
\label{sec:vc2:ingredients}
%──────────────────────────────────────────────────────────────────────────────

\subsection{Kaluza–Klein Spectrum on $S^1_\psi$}
\label{sec:vc2:KK}

The biquaternion field $\Theta(q,\tau)$ with $\tau = t + i\psi$ is quantised on
the imaginary-time circle $S^1_\psi$ (radius $R_\psi$).
The eigenvalues of the kinetic operator $\nabla^\dagger\nabla$ restricted to the
winding sector are:
\begin{equation}
  E_n \;=\; A\,\frac{n^2}{R_\psi^2},
  \qquad n \in \ZZ,
  \label{eq:vc2:KK}
\end{equation}
where $A = \hbar^2/(2m)$ is the standard winding-mode energy coefficient.
Setting natural units $A = 1$ (i.e.\ $R_\psi = 1$, $\hbar = 1$, $2m = 1$)
gives $E_n = n^2$. This is a \emph{choice of units}, not a free parameter.

\begin{resultbox}
\textbf{Source}: \texttt{canonical/appendices/appendix\_alpha\_geometry.tex}
§\ref{sec:modes}, eq.~\eqref{eq:KK_spectrum}.  \Lzero.
\end{resultbox}

\subsection{One-Loop Entropy (Degeneracy) Term}
\label{sec:vc2:entropy}

The one-loop Coleman-Weinberg effective potential receives a contribution
from the functional determinant over quantum fluctuations $\delta\Theta$
around the classical winding background.
In the large-$n$ limit, this determinant is dominated by the logarithm of the
mode density:
\begin{equation}
  S_{\mathrm{CW}}(n)
  \;=\;
  -\frac{1}{2}\ln\det\bigl(-\nabla^\dagger\nabla\bigr)\big|_n
  \;\approx\;
  \frac{N_{\mathrm{eff}}}{2}\,n\ln n
  \;+\; \text{subleading},
  \label{eq:vc2:CW}
\end{equation}
where $N_{\mathrm{eff}} = 12$ is the effective number of charged bosonic modes
(see Theorem~\ref{thm:vc2:Neff}).

The factor $n\ln n$ arises as follows:
\begin{enumerate}
  \item The heat kernel on $S^1_\psi$ at winding number $n$ contributes
        $\ln(n/\Lambda)$ per mode in the UV-regulated functional determinant.
  \item The mode-density factor $n$ multiplies the logarithm because there
        are $n$ distinct winding sectors that contribute coherently to the
        one-loop amplitude at the $n$-th level.
  \item Together: each of the $N_{\mathrm{eff}}$ modes contributes
        $n\ln(n/\Lambda)$; summing and absorbing~$\Lambda$ into the
        renormalisation scheme gives the $B\,n\ln n$ term.
\end{enumerate}

\begin{infobox}
\textbf{Why $n\ln n$ and not $\ln n$}:

The alternative form $\ln n$ (i.e.\ $V = n^2 - B\ln n$, form V4 in the
repository) would arise if the heat-kernel contribution were \emph{independent}
of the winding number — i.e.\ if the functional determinant received the same
logarithmic contribution for every $n$.  This is not the case: the winding
energy $E_n = n^2$ grows with $n^2$, so the one-loop correction must also grow.
The correct large-$n$ behaviour of the Coleman-Weinberg potential is
$\sim n\ln n$, not $\ln n$.

Form V4 is inconsistent with the one-loop calculation and must not be used.
\end{infobox}

\subsection{$N_{\mathrm{eff}} = 12$ from the Biquaternion Algebra}
\label{sec:vc2:Neff}

\begin{theorem}[$N_{\mathrm{eff}} = 12$, \PROVED\ \Lzero]
\label{thm:vc2:Neff}
The biquaternion algebra $\BB = \CC\otimes_\RR\HH$ admits exactly
$N_{\mathrm{eff}} = 12$ independent real bosonic modes contributing to the
one-loop vacuum polarisation:
\[
  N_{\mathrm{eff}}
  \;=\;
  N_{\mathrm{phases}} \times N_{\mathrm{helicity}} \times N_{\mathrm{charge}}
  \;=\;
  3 \times 2 \times 2
  \;=\; 12,
\]
with $N_{\mathrm{phases}} = \dim_\RR(\mathrm{Im}\,\HH) = 3$.
\end{theorem}

\noindent\textbf{Source}: \texttt{canonical/n\_eff/step3\_N\_eff\_result.tex}.

%──────────────────────────────────────────────────────────────────────────────
\section{Derivation of $V_{\mathrm{eff}}(n)$}
\label{sec:vc2:derivation}
%──────────────────────────────────────────────────────────────────────────────

\begin{theorem}[One-loop effective potential, \PROVED\ \Lone]
\label{thm:vc2:Veff}
The one-loop effective potential for the winding mode $n$ on $S^1_\psi$,
derived from the Coleman-Weinberg functional determinant of
$\nabla^\dagger\nabla$ on the $\BB$-valued field $\Theta$, is:
\begin{equation}
  \boxed{V_{\mathrm{eff}}(n) \;=\; n^2 \;-\; B\,n\ln n,}
  \qquad n \in \ZZ_{>0},
  \label{eq:vc2:Veff}
\end{equation}
where the coefficient $B$ collects the one-loop vacuum-polarisation contributions
of all $N_{\mathrm{eff}} = 12$ charged modes.
\end{theorem}

\begin{proof}
Combine the classical energy $E_n = n^2$ (natural units, §\ref{sec:vc2:KK})
with the one-loop entropy term $S_{\mathrm{CW}}(n) = \frac{N_{\mathrm{eff}}}{2}\,n\ln n$
(eq.~\eqref{eq:vc2:CW}):
\[
  V_{\mathrm{eff}}(n)
  \;=\;
  E_n - S_{\mathrm{CW}}(n)
  \;=\;
  n^2 - B\,n\ln n,
  \qquad
  B \;\equiv\; \frac{N_{\mathrm{eff}}}{2} = 6
  \quad (\text{baseline from entropy alone}).
\]
The full one-loop coefficient receives additional contributions from the
functional determinant over the imaginary-time modes; the complete one-loop
baseline is $B_0 = 8\pi \approx 25.13$ (Theorem~\ref{thm:vc2:B0}), derived
from the vacuum polarisation integral with $N_{\mathrm{eff}} = 12$ modes.
The functional form $n^2 - B\,n\ln n$ is fixed by the calculation;
the open problem is the precise value of $B$.
$\square$
\end{proof}

\begin{theorem}[Baseline coefficient, \PROVED\ \Lone]
\label{thm:vc2:B0}
The one-loop vacuum-polarisation contribution from $N_{\mathrm{eff}} = 12$
charged modes on $S^1_\psi$ gives:
\[
  B_0 \;=\; \frac{2\pi N_{\mathrm{eff}}}{3} \;=\; 8\pi \;\approx\; 25.133.
\]
\end{theorem}

\noindent\textbf{Source}:
\texttt{canonical/n\_eff/step2\_vacuum\_polarization.tex}, Theorem~3.1.

%──────────────────────────────────────────────────────────────────────────────
\section{Stationarity Analysis}
\label{sec:vc2:stationary}
%──────────────────────────────────────────────────────────────────────────────

\begin{theorem}[Stationarity condition, \PROVED\ \Lone]
\label{thm:vc2:stat}
The continuous minimum of $V_{\mathrm{eff}}(n) = n^2 - B\,n\ln n$
satisfies the transcendental equation:
\begin{equation}
  \frac{\partial V_{\mathrm{eff}}}{\partial n}\bigg|_{n^*}
  \;=\;
  2n^* - B\bigl(\ln n^* + 1\bigr)
  \;=\; 0,
  \label{eq:vc2:dV}
\end{equation}
equivalently
\begin{equation}
  \boxed{2n^* \;=\; B\,\bigl(\ln n^* + 1\bigr).}
  \label{eq:vc2:stat}
\end{equation}
This equation has \textbf{no elementary closed form}.
\end{theorem}

\begin{proof}
Direct differentiation:
$\partial/\partial n\,(n^2) = 2n$;
$\partial/\partial n\,(B\,n\ln n) = B\,(\ln n + 1)$.
Setting the derivative to zero yields eq.~\eqref{eq:vc2:stat}.
The function $f(n) = 2n - B(\ln n + 1)$ is strictly increasing for $n > B/2$
and strictly decreasing for $n < B/2$, so the solution is unique and a minimum.
$\square$
\end{proof}

\begin{corollary}[Second derivative, \PROVED]
\label{cor:vc2:2nd}
\[
  \frac{\partial^2 V_{\mathrm{eff}}}{\partial n^2}\bigg|_{n^*}
  = 2 - \frac{B}{n^*}
  > 0
  \;\Longleftrightarrow\;
  n^* > \frac{B}{2}.
\]
For $B \approx 46.3$ and $n^* = 137$: $137 > 23.15$ \,\checkmark.
The stationary point is a \textbf{minimum}.
\end{corollary}

\begin{table}[h]
\centering
\caption{Numerical solutions of eq.~\eqref{eq:vc2:stat} for selected $B$ values}
\label{tab:vc2:nstar}
\begin{tabular}{llll}
\toprule
$B$ value & $n^*_{\mathrm{continuous}}$ & Prime minimiser & Note \\
\midrule
$B_0 = 8\pi \approx 25.133$ & $65.0$ & $p = 67$ (or 61?) & proved [L1] \\
$B_{\mathrm{base}} = N_{\mathrm{eff}}^{3/2} \approx 41.569$ & $120.4$ & $p = 127$ & conjectural [MC] \\
$\mu(\Gamma_0(137))/3 = 46.000$ & $136.0$ & $p = 137$ & exact math; 0.64\% below $B_{\mathrm{phenom}}$ \\
$B_{\mathrm{required}} = 46.284$ & $137.000$ & $p = 137$ & exact constraint \\
$B_{\mathrm{phenom}} \approx 46.298$ & $137.05$ & $p = 137$ & phenomenological \\
\bottomrule
\end{tabular}
\end{table}

\begin{proposition}[Required $B$ for $n^* = 137$, \PROVED]
\label{prop:vc2:Breq}
Substituting $n^* = 137$ into eq.~\eqref{eq:vc2:stat}:
\begin{equation}
  B_{\mathrm{required}}
  \;=\;
  \frac{2 \times 137}{\ln 137 + 1}
  \;=\;
  \frac{274}{5.91998\ldots}
  \;\approx\; 46.284.
  \label{eq:vc2:Breq}
\end{equation}
The phenomenological value $B_{\mathrm{phenom}} = 46.298$ in the repository
differs from $B_{\mathrm{required}}$ by $0.030\%$ — consistent with rounding in
earlier documents.
\end{proposition}

%──────────────────────────────────────────────────────────────────────────────
\section{Derivation Table for All V_eff Variants}
\label{sec:vc2:variants}
%──────────────────────────────────────────────────────────────────────────────

The following table summarises all $V_{\mathrm{eff}}$ variants found in the
repository, their derivatives, stationarity conditions, and the resulting $n^*$.

\begin{center}
\renewcommand{\arraystretch}{1.4}
\begin{tabular}{p{3.5cm} p{3cm} p{3.5cm} p{2.5cm} p{2.5cm}}
\toprule
\textbf{Form} & \textbf{$\partial V/\partial n$} & \textbf{Stationarity}
  & \textbf{$n^*$ (cont.)} & \textbf{$B$ for $n^*=137$} \\
\midrule
V1: $n^2 - B\,n\ln n$ \newline \textbf{(correct)}
  & $2n - B(\ln n + 1)$
  & $2n^* = B(\ln n^*+1)$ \newline transcendental
  & $137.0$ for $B\approx 46.28$
  & $B = 274/(\ln 137+1) \approx 46.284$ \\
\midrule
V4: $n^2 - B\ln n$ \newline \textbf{(wrong — V4)}
  & $2n - B/n$
  & $n^* = \sqrt{B/2}$ \newline closed form
  & $\approx 4.8$ for $B\approx 46.3$
  & $B = 2\times 137^2 = 37538$ \\
\bottomrule
\end{tabular}
\end{center}

\begin{warnbox}
Only form V1 gives $n^* = 137$ for a physically reasonable coefficient
$B \approx 46.3$.

Form V4 gives $n^* \approx 4.8$ for the same $B$, which is nowhere near 137.
The stationarity formula $n^* = \sqrt{B/2}$ that appears in older repository
documents belongs to form V4 and is a \textbf{transcription error}.
\end{warnbox}

%──────────────────────────────────────────────────────────────────────────────
\section{Classification of All $B$ Values}
\label{sec:vc2:Bclass}
%──────────────────────────────────────────────────────────────────────────────

\begin{center}
\renewcommand{\arraystretch}{1.3}
\begin{tabular}{p{4cm} p{1.8cm} p{3.5cm} p{3.5cm}}
\toprule
\textbf{Symbol} & \textbf{Value} & \textbf{Classification} & \textbf{Proof status} \\
\midrule
$B_0 = 8\pi$ & $\approx 25.13$ & DERIVED & \PROVED\ \Lone \\
$B_{\mathrm{base}} = N_{\mathrm{eff}}^{3/2}$ & $\approx 41.57$ & CONJECTURAL & \MC; Gap G3-k open \\
$\mu(\Gamma_0(137))/3$ & $46.00$ & EXACT MATH & \PROVED\ \Lzero\ (index formula); ID with $B$ is \OPEN \\
$B_{\mathrm{required}}$ & $46.284$ & EXACT CONSTRAINT & Derived from $n^*=137$; circular if used as input \\
$B_{\mathrm{phenom}}$ & $\approx 46.298$ & PHENOMENOLOGICAL & Back-solved from $n^*=137$; \PHENOM \\
$B_{\mathrm{full}} = B_{\mathrm{base}}\cdot R$ & $\approx 46.31$ & OPEN & $R \approx 1.114$; no derivation of $R$ \\
$B_{\mathrm{obs}} = N_{\mathrm{eff}}^{3/2}(2\eta(i))^{c/12}$ & $\approx 46.281$ & OBSERVATION & $0.007\%$ from $B_{\mathrm{required}}$; \OBS\ pending proof \\
\bottomrule
\end{tabular}
\end{center}

\begin{gapbox}
\textbf{Gap G137-B (primary blocker)}: Derive $B \approx 46.284$ from
$S[\Theta]$ without using $\alpha$ or $n^*=137$ as input.

The missing normalisation factor is now identified as
$R = (2\eta(i))^{c/12} = (2\eta(i))^{1/4} = (\Gamma(1/4)/\pi^{3/4})^{1/4}
\approx 1.1132$, where $c = 3$ is the central charge of the UBT CFT and
$\eta(i) = \Gamma(1/4)/(2\pi^{3/4})$ is the Dedekind eta at the self-dual point
$\tau = i$.

Formula: $B_{\mathrm{obs}} = N_{\mathrm{eff}}^{3/2}(2\eta(i))^{c/12} \approx 46.281$
matches $B_{\mathrm{required}} = 46.284$ to $0.007\%$ with no free parameters.
Proof of this identification is required to close Gap G137-B.

See \texttt{research\_tracks/T3\_ALPHA/cw\_determinant\_full\_derivation.tex}
for the full derivation and evidence.
\end{gapbox}

%──────────────────────────────────────────────────────────────────────────────
\section{Errors Corrected by This Document}
\label{sec:vc2:errors}
%──────────────────────────────────────────────────────────────────────────────

\begin{center}
\renewcommand{\arraystretch}{1.2}
\begin{tabular}{p{4.5cm} p{2.5cm} p{4.5cm} p{2cm}}
\toprule
\textbf{Document} & \textbf{Section/eq} & \textbf{Error} & \textbf{Severity} \\
\midrule
\texttt{appendix\_alpha\_geometry.tex} (pre-2026-05-02)
  & \S\ref{sec:Veff}, eq.\ \eqref{eq:Veff}
  & $V = n^2 - B\ln n$ (missing $n$); leads to wrong $n^* \approx 4.8$
  & HIGH \\
\texttt{appendix\_alpha\_geometry.tex} (pre-2026-05-02)
  & \S\ref{sec:prime}, Thm., proof
  & $n^* = \sqrt{B/2}$ stated as stationarity of wrong V4 form
  & HIGH \\
\texttt{ALPHA\_PROGRESS\_REPORT.md} (pre-2026-05-02)
  & \S2.3 header, \S9 entry A7
  & V\_eff form and clean-chain entry showed wrong V4 variant
  & MEDIUM \\
\texttt{alpha\_equation\_matrix.tex} (pre-2026-04-29)
  & Thm.~[L1], shared foundation
  & $n^*(B) = e^{(2-B)/B}\cdot e$ (equals $e^{2/B}\approx 1.04$, not 137)
  & HIGH (fixed in prior session) \\
\bottomrule
\end{tabular}
\end{center}

%──────────────────────────────────────────────────────────────────────────────
\section{Summary and Next Steps}
\label{sec:vc2:summary}
%──────────────────────────────────────────────────────────────────────────────

\begin{tcolorbox}[colback=blue!4!white,colframe=blue!50!black,
                  title=Summary of Findings]
\begin{enumerate}
  \item \textbf{The correct form of $V_{\mathrm{eff}}$ is $n^2 - B\,n\ln n$} (V1),
        derived from the one-loop Coleman-Weinberg functional determinant on
        $S^1_\psi$.  The alternative form $n^2 - B\ln n$ (V4) is inconsistent
        with the one-loop calculation.

  \item \textbf{The stationarity condition is transcendental}: $2n^* = B(\ln n^*+1)$.
        There is no closed-form solution.  The formula $n^* = \sqrt{B/2}$
        belongs to the wrong form V4 and has been removed from all relevant
        documents.

  \item \textbf{$n^* = 137$ is a valid prime minimiser} of the correct $V_{\mathrm{eff}}$
        for $B \in [46.00, 46.50]$.  The exact constraint is $B_{\mathrm{required}}
        \approx 46.284$.

  \item \textbf{$B \approx 46.3$ is phenomenological.}  The only derived
        value is $B_0 = 8\pi \approx 25.13$.  The gap is a factor
        $\approx 1.84$ that must be closed without using $\alpha$ as input
        (Gap G137-B).

  \item \textbf{The $\alpha$ route is mathematically consistent but not yet complete.}
        It is valid as a conditional result (given $B \approx 46.3$) and
        must not be promoted to a first-principles proof until Gap G137-B is
        closed.
\end{enumerate}
\end{tcolorbox}

\paragraph{Next steps.}
\begin{enumerate}
  \item Attack Gap G137-B via modular bootstrap on $\hat{Z}(\tau) = \vartheta_3^3(\tau)$.
  \item Explore whether $\mu(\Gamma_0(137))/3 = 46.00$ can be derived as $B$
        from the Hecke eigenform structure of $S[\Theta]$
        (see \texttt{canonical/alpha/alpha\_equation\_matrix.tex}, Route~A).
  \item Do not promote $n^* = 137$ to PROVED until Gap G137-B is closed.
\end{enumerate}

\ifdefined\INCLUDEMODE
\else
  \end{document}
\fi
